% !TEX TS-program = XeLaTeX
% use the following command:
% all document files must be coded in UTF-8
\documentclass[portuguese]{textolivre}
% build HTML with: make4ht -e build.lua -c textolivre.cfg -x -u article "fn-in,svg,pic-align"

\journalname{Texto Livre}
\thevolume{19}
%\thenumber{1} % old template
\theyear{2026}
\receiveddate{\DTMdisplaydate{2025}{11}{7}{-1}} % YYYY MM DD
\accepteddate{\DTMdisplaydate{2026}{1}{5}{-1}}
\publisheddate{\DTMdisplaydate{2026}{6}{16}{-1}}
\corrauthor{Hugo de Oliveira}
\articledoi{10.1590/1983-3652.2026.62731}
%\articleid{NNNN} % if the article ID is not the last 5 numbers of its DOI, provide it using \articleid{} commmand 
% list of available sesscions in the journal: articles, dossier, reports, essays, reviews, interviews, editorial
\articlesessionname{articles}
\runningauthor{Oliveira e Pinto} 
%\editorname{Leonardo Araújo} % old template
\sectioneditorname{Daniervelin Pereira~\orcid{0000-0003-1861-3609}}
\layouteditorname{Leonardo Araujo~\orcid{0000-0003-3884-2177}}

\title{Capacidade de agir em infraestruturas digitais: pode a performatividade contribuir para a compreensão da subjetivação tecnológica?}
\othertitle{Agency in digital infrastructures: can performativity contribute to the understanding of technological subjectivation?}
% if there is a third language title, add here:
%\othertitle{Artikelvorlage zur Einreichung beim Texto Livre Journal}

\author[1]{Hugo de Oliveira~\orcid{0000-0001-6412-8116}\thanks{Email: \href{mailto:hugodeoliveira03@gmail.com}{hugodeoliveira03@gmail.com}}}
\author[2]{Joana Plaza Pinto~\orcid{0000-0001-8052-9390}\thanks{Email: \href{mailto:joplazapinto@ufg.br}{joplazapinto@ufg.br}}}
\affil[1]{Universidade Federal de Goiás, Faculdade de Ciências Sociais, Programa de Pós-Graduação em Sociologia, Goiânia, GO, Brasil.}
\affil[2]{Universidade Federal de Goiás, Faculdade de Letras, Departamento de Linguística e Língua Portuguesa,
Goiânia, GO, Brasil.}

\addbibresource{article.bib}
% use biber instead of bibtex
% $ biber article

% used to create dummy text for the template file
\definecolor{dark-gray}{gray}{0.35} % color used to display dummy texts
\usepackage{lipsum}
\SetLipsumParListSurrounders{\colorlet{oldcolor}{.}\color{dark-gray}}{\color{oldcolor}}

% used here only to provide the XeLaTeX and BibTeX logos
\usepackage{hologo}

% if you use multirows in a table, include the multirow package
\usepackage{multirow}

% provides sidewaysfigure environment
\usepackage{rotating}

% CUSTOM EPIGRAPH - BEGIN 
%%% https://tex.stackexchange.com/questions/193178/specific-epigraph-style
\usepackage{epigraph}
\renewcommand\textflush{flushright}
\makeatletter
\newlength\epitextskip
\pretocmd{\@epitext}{\em}{}{}
\apptocmd{\@epitext}{\em}{}{}
\patchcmd{\epigraph}{\@epitext{#1}\\}{\@epitext{#1}\\[\epitextskip]}{}{}
\makeatother
\setlength\epigraphrule{0pt}
\setlength\epitextskip{0.5ex}
\setlength\epigraphwidth{.7\textwidth}
% CUSTOM EPIGRAPH - END

% to use IPA symbols in unicode add
%\usepackage{fontspec}
%\newfontfamily\ipafont{CMU Serif}
%\newcommand{\ipa}[1]{{\ipafont #1}}
% and in the text you may use the \ipa{...} command passing the symbols in unicode

% LANGUAGE - BEGIN
% ARABIC
% for languages that use special fonts, you must provide the typeface that will be used
% \setotherlanguage{arabic}
% \newfontfamily\arabicfont[Script=Arabic]{Amiri}
% \newfontfamily\arabicfontsf[Script=Arabic]{Amiri}
% \newfontfamily\arabicfonttt[Script=Arabic]{Amiri}
%
% in the article, to add arabic text use: \textlang{arabic}{ ... }
%
% RUSSIAN
% for russian text we also need to define fonts with support for Cyrillic script
% \usepackage{fontspec}
% \setotherlanguage{russian}
% \newfontfamily\cyrillicfont{Times New Roman}
% \newfontfamily\cyrillicfontsf{Times New Roman}[Script=Cyrillic]
% \newfontfamily\cyrillicfonttt{Times New Roman}[Script=Cyrillic]
%
% in the text use \begin{russian} ... \end{russian}
% LANGUAGE - END

% EMOJIS - BEGIN
% to use emoticons in your manuscript
% https://stackoverflow.com/questions/190145/how-to-insert-emoticons-in-latex/57076064
% using font Symbola, which has full support
% the font may be downloaded at:
% https://dn-works.com/ufas/
% add to preamble:
% \newfontfamily\Symbola{Symbola}
% in the text use:
% {\Symbola }
% EMOJIS - END

% LABEL REFERENCE TO DESCRIPTIVE LIST - BEGIN
% reference itens in a descriptive list using their labels instead of numbers
% insert the code below in the preambule:
%\makeatletter
%\let\orgdescriptionlabel\descriptionlabel
%\renewcommand*{\descriptionlabel}[1]{%
%  \let\orglabel\label
%  \let\label\@gobble
%  \phantomsection
%  \edef\@currentlabel{#1\unskip}%
%  \let\label\orglabel
%  \orgdescriptionlabel{#1}%
%}
%\makeatother
%
% in your document, use as illustraded here:
%\begin{description}
%  \item[first\label{itm1}] this is only an example;
%  % ...  add more items
%\end{description}
% LABEL REFERENCE TO DESCRIPTIVE LIST - END


% add line numbers for submission
%\usepackage{lineno}
%\linenumbers

\begin{document}
\maketitle

\begin{polyabstract}
\begin{abstract}
Este artigo visa discutir os modos como concebemos a correlação entre linguagem, subjetivação e infraestruturas técnicas digitais a partir da noção de performativo de J. L. Austin. Para isso, nosso ponto de partida é a análise de Fanon sobre as transformações subjetivas provocadas pelo uso do rádio durante a Revolução Argelina comparadas ao processo contemporâneo de digitalização da vida analisado por Letícia Cesarino. Tais estudos evidenciam, à sua maneira, o caráter performativo da linguagem em criar realidades. Contudo, diferentes apropriações da obra de Austin trazem contribuições divergentes para os estudos da linguagem nas sociedades digitais: por um lado, a normatização do performativo e a oposição verdadeiro x falso, presentes em Benveniste, Searle e Cesarino; por outro lado, a crítica focada nos efeitos da linguagem no mundo presente em Derrida e Butler. Analisando as possibilidades de uso da “verdade” na compreensão da linguagem em ambientes virtuais, posicionamo-nos na vertente crítica, afinada à nossa leitura de Fanon, para a compreensão das tecnologias nos processos de subjetivação como atos de fala performativos, nos quais a capacidade de agir não é eliminada, mas reconfigurada como agência situada.

\keywords{Linguagem \sep Performatividade \sep Subjetivação \sep Infraestruturas digitais \sep Agência}
\end{abstract}

\begin{english}
\begin{abstract}
This article aims to discuss how we understand the correlation between language, subjectivation, and digital technical infrastructures based on J. L. Austin’s notion of the performative. For that, our stating point is Fanon’s analysis of the subjective transformations provoked by the use of radio during the Algerian Revolution, comparing it to the contemporary process of digitalization of life as analyzed by Letícia Cesarino. Each of these studies, highlight the performative characteristic of language in creating realities rather than merely describing them. However, different appropriations of Austin’s work bring divergent contributions to the study of language in digital societies: on the one hand, the normalization of the performative and the opposition between true and false, present in Benveniste, Searle, and Cesarino; on the other, the critique focused on the effects of language on the world present in Derrida and Butler. By analyzing the possibilities of using “truth” in the understanding of language in digital environments, we position ourselves in the critical perspective, in line with our reading of Fanon, to understand technologies in the processes of subjectivation as performative speech acts, in which the agency is not eliminated, but reconfigured as a situated one.

\keywords{Language \sep Performativity \sep Subjectivation \sep Digital infrastructures \sep Agency}
\end{abstract}
\end{english}
% if there is another abstract, insert it here using the same scheme
\end{polyabstract}

\section{Introdução}\label{sec-intro}

\epigraph{“Ter um aparelho de rádio é, solenemente, entrar na guerra.”}{Frantz Fanon}


Durante a Guerra de Libertação da Argélia (1954-1962), uma nova tecnologia, o rádio, foi fundamental para o desfecho do conflito e a conquista da independência em relação à França. O papel dessa tecnologia no combate foi ambivalente: por um lado, na medida em que transmitia os valores da sociedade colonial, representou um símbolo da presença francesa; por outro lado, na medida em que potencializava uma vasta rede de informações, o rádio produziu processos de subjetivação que impulsionaram a Revolução Argelina. Esse evento complexo foi analisado pelo psiquiatra e filósofo martinicano \textcite{fanon2021argelia}, num artigo originalmente publicado em 1959, no qual ele aponta as transformações subjetivas engendradas por mudanças tecnológicas, e indica, sem assim o nomear, o caráter performativo da tecnologia: “À escuta da Revolução [no rádio], o argelino existe com ela, a faz existir” \cite[p. 692]{fanon2021argelia}.

Inicialmente, durante a dominação francesa, o rádio tinha uma presença restrita entre os colonos. Em um contexto de estratificação econômica entre as sociedades dominante e dominada, a rádio representava uma maneira de “manter-se civilizado”, ela lembrava os dominadores de um poder e o distribuía entre eles: “[...] a rádio é, em geral, vivida como laço com o mundo civilizado, como instrumento de eficácia de resistência à influência corrosiva de uma sociedade indígena imóvel, sem perspectiva, atrasada e sem valor” \cite[p. 679]{fanon2021argelia}. Diante do fechamento desse mundo colonialista em si, o povo argelino pouco se interessou por um aparelho que representava, sobretudo, o mundo colonial falando. Até o final da década de 1940, o rádio enfrentou uma indiferença e um desinteresse muito fortes porque não havia uma relação de “verdade” entre a informação veiculada por esse aparelho e as condições políticas e sociais da população argelina. Contudo, tal situação se transformou com a criação de emissoras nacionais na Síria, no Egito e no Líbano nos últimos anos de 1940, e com a união da Argélia, em 1951, à Frente do Magreb anticolonialista.

Segundo \textcite{fanon2021argelia}, somada a isso, a percepção de que algo maior estava em andamento levou os argelinos a buscarem e disporem de suas próprias informações, com a finalidade de autodefesa e de escapar às “verdades” dos opressores.

\begin{quote}
    Saber o que se passa, conhecer ao mesmo tempo as perdas reais do inimigo e as próprias, torna-se fundamental. O argelino, nesta época, tem a necessidade de içar sua vida ao nível da Revolução. Ele tem a necessidade de entrar na vasta rede de informação; ele tem a necessidade de se introduzir em um mundo onde as coisas se passam, onde o evento existe, onde as forças agem. O argelino, através da existência de uma guerra feita pelo seu povo, leva uma comunidade ao ato. Às informações do inimigo, o argelino deve opor suas próprias informações. À verdade do opressor, outrora rejeitada como mentira absoluta, é contraposta, enfim, outra verdade concreta \cite[p. 682]{fanon2021argelia}.
\end{quote}

Tal processo se intensificou com o início das transmissões da rádio Voix de l’Algérie Libre (Voz da Argélia Livre), que adquiriu um valor essencial por enunciar a mensagem da Revolução. A partir desse momento, o aparelho de rádio deixou de ser uma forma de opressão cultural e tornou-se um meio para resistir ao opressor: “[...] a sociedade argelina, por um movimento autônomo interno, decide assumir a nova técnica e estar, assim, ligada aos novos sistemas de sinalização postos no mundo pela Revolução” \cite[p. 687]{fanon2021argelia}. Fanon argumenta que a apropriação do rádio levou à contestação do princípio da dominação estrangeira, provocando mutações essenciais na consciência dos colonizados e na sua situação de ser humano perante o mundo.

A análise de Fanon aponta para a percepção de um mundo cindido, no qual há uma divisão profunda entre colonizadores e colonizados. Nele, a violência indiscriminada se desdobra na pilhagem de esquemas culturais, engendrando posições sociais hierarquizadas via marcas fenotípicas e culturais e divisão racial do trabalho. Nesse contexto, o racismo “é tanto um ‘produto’ quanto um processo pelo qual o grupo dominante lança mão para desarticular as possíveis linhas de força do dominado” \cite[p. 153]{faustino2018fanon}.

No reconhecimento dos impactos das tecnologias de comunicação, é inevitável ressaltar o papel da linguagem na imposição da ordem colonial. Como argumenta \textcite[p. 680]{fanon2021argelia}:

\begin{quote}
    Num outro nível, como sistema de informação, como portadora da linguagem, portanto da mensagem, a estação de Rádio pode ser apreendida no seio da situação colonial de maneira particular. A tecnologia radiofônica, a imprensa e, de uma maneira geral, os sistemas, as mensagens, os transmissores de signos, existem na sociedade colonial segundo um estatuto perfeitamente diferenciado. A sociedade argelina, a sociedade dominada nunca participa deste mundo de signos.
\end{quote}

Um conjunto de discursos baseados na inferiorização e na racialização da população subalternizada se enunciou, frequentemente, via “um mundo de signos” que se afirmava constativo \cite{austin1962how}, ou seja, supostamente ou “verdadeiro” ou “falso”, descrevendo ou relatando um estado de coisas. Contudo, restringindo-nos aos termos “verdadeiro” ou “falso”, pouco entenderíamos o papel da linguagem no jogo de “forças”, do “ato”, da contraposição de uma “outra verdade concreta” contra o sistema colonial. Afinal, mais que descrever as ações, a linguagem realiza ações no mundo \cite[p. 6]{austin1962how}.

Nesse sentido, quando enunciava os discursos franceses sobre a superioridade civilizatória da colonização, a rádio criava uma realidade na qual os nativos não passam de seres atrasados salvos pelo mito da modernidade. Em contrapartida, é por meio da mesma tecnologia que também se cria uma realidade na qual a independência e a Revolução se tornam possíveis por meio da formação da voz argelina.

Embora o texto de Fanon tenha sido publicado ainda em 1959, é possível traçar paralelos entre aquele contexto e o atual, no qual transformações tecnológicas têm provocado profundas transformações políticas e sociais. Com o avanço das ciências, do mercado da informação e, consequentemente, do chamado Capitalismo 4.0 ou digital, aparelhos tecnológicos como o computador e os smartphones tornaram-se ainda mais fundamentais para os modos de subjetivação. Esses processos sãos mediados pela linguagem: os sujeitos se constituem tanto por aquilo que leem, falam, ouvem, reagem e digitam, quanto pelo meio no qual fazem essas ações linguísticas.

Nos últimos anos, as novas tecnologias de informação – impulsionadas por grandes empresas localizadas, em sua maioria, no Vale do Silício nos Estados Unidos – têm provocado uma profunda transformação da esfera pública, marcada pela progressiva digitalização da vida e da política. Nesse contexto, no Brasil e no mundo, uma profunda reorganização do campo político-identitário foi provocada pelo aumento da capilaridade de infraestruturas técnicas, como as plataformas digitais. Assim, a partir de redes sociais, de algoritmos, do compartilhamento e da extração de dados e da reformulação de processos de mediação online têm emergido novos modelos interacionais \cite{pinto2024novaordem} e novas dinâmicas de formação de subjetividades ancoradas na cibernética, que constituem uma nova economia subjetiva \cite{cesarino2022b}.

De forma similar ao mundo colonial analisado por Fanon, essa nova realidade tecno-social se apresenta como um mundo cindido, no qual a polarização política é a base de uma divisão amigo-inimigo que coloca o sistema de peritos das democracias liberais em cheque. Algumas pesquisas chegam a argumentar que, a partir de uma lógica algorítmica, desenvolve-se uma eu-pistemologia, um sistema de referências que leva em consideração exclusivamente o indivíduo, constituído por dinâmicas de clusterização e feeds personalizados, que provoca o colapso de instituições sociais tradicionalmente aceitas, como a democracia, a ciência e a “verdade” \cite{cesarino2021b}.

Diante desse cenário, nos perguntamos como fica a linguagem, “o mundo de signos” que interrogou Fanon, e como o performativo, a força do ato na linguagem, poderia nos ajudar a entender os impactos das infraestruturas digitais no mundo contemporâneo. Nosso objetivo neste artigo é debater o lugar da linguagem nos processos de subjetivação mediados pela tecnologia, seguindo a noção de performatividade nos estudos sobre ambientes digitais e propondo uma elaboração crítica dessa noção diante desses ambientes.

Para isso, partimos dos riscos e da capacidade criativa, dos limites e do potencial da linguagem a partir de suas características performativas retomando os estudos do filósofo inglês John L. Austin em diálogo com as reflexões de Fanon sobre as tecnologias no colonialismo e no movimento anticolonialista.

Na seção seguinte, fazemos uma revisão da Teoria dos Atos de Fala e debates decorrentes envolvendo Émile Benveniste, John Searle e Kanavillil Rajagopalan. Em seguida, trazemos alguns debates sobre a articulação de linguagem, subjetivação e infraestruturas digitais, com destaque para a apropriação e o percurso dos conceitos de performativo e de performatividade nos estudos da antropóloga Letícia Cesarino.

Na penúltima seção, apresentamos uma alternativa crítica à visão que identificamos no percurso de Cesarino a partir das apropriações críticas da Teoria da Performatividade formuladas por Jacques Derrida e Judith Butler, tendo como pano de fundo nossa leitura de Fanon e intérpretes sobre as tecnologias no processo de subjetivação. Finalizamos rearticulando este debate com o potencial do performativo para o avanço da discussão sobre a conexão entre linguagem, subjetivação e infraestruturas digitais.

\section{Uma visão performativa da linguagem}\label{sec-normas}
Segundo \textcite{ottoni2002austin}, a filosofia de Austin surgiu em um momento histórico preciso. Ocorria o fortalecimento das discussões sobre linguagem na Escola de Oxford na década de 1940, a gênese da sintaxe gerativa na Linguística estadunidense na década de 1950 com Chomsky, o desenvolvimento de uma perspectiva linguística centrada na enunciação com Benveniste na França, além da elaboração de diversas reflexões sobre a linguagem por autores como Foucault, Lacan e Derrida, que desembocariam numa perspectiva pós-estruturalista. Nesse contexto, a obra de Austin é reconhecida por ter revolucionado a Filosofia Analítica por meio do estudo da linguagem ordinária, principalmente no que diz respeito à questão da verdade, pano de fundo para a distinção fundamental entre enunciados performativos e constativos. Além disso, seu trabalho é reconhecido por não respeitar fronteiras disciplinares, pois, para ele, não era possível pensar a linguagem de forma compartimentada, o que levou à quebra de barreiras entre a Filosofia e a Linguística \cite{ottoni2002austin}.

Suas reflexões sobre constativo-performativo partem da discussão sobre a verdade, a partir de uma crítica à falácia descritiva, a linguagem não é puramente descritiva, uma vez que existem situações nas quais não descrevemos ações, mas as praticamos. Inicialmente, o filósofo inglês afirma que os enunciados constativos, evocados sob o nome de afirmação pelos filósofos, carregariam a propriedade de serem verdadeiros ou falsos, enquanto os performativos, por sua vez, não poderiam ser nem um nem outro, uma vez que têm sua própria função, realizar uma ação \cite[p. 5–6]{austin1962how}. Segundo essa perspectiva, há uma questão de referência que envolve a relação linguagem-mundo, enquanto os primeiros evocariam uma referência supostamente objetiva, isso não ocorreria no segundo grupo, pois eles realizam uma ação.

Assim, diferentemente do constativo que seria verdadeiro ou falso, o enunciado performativo pode ser feliz, se realizar a ação pretendida, ou infeliz, se a ação não se realizar. Algumas condições de infelicidade específicas dos atos de fala foram descritas por \textcite{austin1998performativo}: a) a nulidade, quando quem enuncia não está em condições de realizar o ato de fala, como dizer “você herdará meu relógio” quando eu não tenho relógio; b) o abuso de fórmula ou a falta de sinceridade, quando quem enuncia não se compromete com o ato ligado à enunciação, como dizer “eu prometo” sem a intenção de cumprir a promessa; c) a quebra de compromisso, quando quem enuncia não se compromete com os atos posteriores à enunciação, como dizer “seja bem-vindo” e em seguida tratar o interlocutor como um estranho.

A princípio, o filósofo chega a estabelecer duas formas normais de expressão do performativo: a primeira na qual há um verbo no início do enunciado na primeira pessoa do singular, no presente do indicativo, na voz ativa, por exemplo, “eu prometo que...”; na segunda, há um verbo na voz passiva, na segunda ou na terceira pessoa do presente do indicativo, por exemplo, “os passageiros estão convidados a amarrar os cintos”. Contudo, ele chega à conclusão de que há enunciados performativos que não se expressam por essas formas, pois os atos ilocucionários permitem a existência de enunciados performativos sem uma forma gramatical pré-estabelecida \cite[p. 53–66]{austin1962how}. Dessa forma, é possível afirmar que, em cada diferente tipo de enunciado, há um performativo mascarado ou performativos implícitos \cite{ottoni2002austin}. Essa argumentação leva a consequências epistemológicas radicais nos estudos sobre a linguagem, uma vez que:

\begin{quote}
    Este salto, que desfaz a distinção entre performativo-constativo produz uma visão de linguagem que não é mais idêntica à utilizada na distinção anterior entre performativo e o constativo. Esta visão produz, uma virada brutal na questão da referência; ou seja, verdade e falsidade são conceitos que não terão mais um papel relevante nem prioritário para Austin. A partir deste momento podemos falar de uma visão performativa, na qual o sujeito não pode se desvincular de seu objeto de fala e, consequentemente, não é possível analisar este objeto fala desvinculado do sujeito \cite[p. 130]{ottoni2002austin}.
\end{quote}

Contudo, os desdobramentos e contrapontos da teoria produzida por Austin nem sempre destacam essa visão, reforçando uma distinção entre Filosofia e Linguística e enfatizando, inclusive, uma posição epistemológica que remonta à normatividade das formas linguísticas em vez de seus efeitos no mundo. Esse é o caso proeminente do posicionamento de \textcite{benveniste1991filosofia} acerca dos enunciados performativos. Para o linguista francês, o fenômeno do performativo levaria a um bom relacionamento entre as duas áreas, desde que cada uma delimitasse muito bem a sua posição diante da linguagem, uma visão contrária aos questionamentos anti-disciplinares de Austin.

Para Benveniste, é necessário manter a distinção performativo-constativo, uma vez que há critérios formais que legitimam essa divisão. Retomando as definições formais primárias do performativo, ele descreve uma série de condições nas quais esses enunciados podem ser descritos, chegando à conclusão de que:

\begin{quote}
    Não vemos, portanto, razão para abandonar a distinção entre performativo e constativo. Acreditamo-la justificada e necessária, com a condição de que a mantenhamos dentro das condições estritas de emprego que a autorizam, sem fazer intervir a consideração do ‘resultado obtido’ que é fonte de confusão \cite[p. 305]{benveniste1991filosofia}.
\end{quote}

A partir disso, é possível observar que o autor permanece focado no nível linguístico, utilizando o performativo de uma maneira estrita \cite{ottoni2002austin}. Com isso, ele reforça uma visão epistemológica da linguagem na qual seus efeitos não fazem parte do problema linguístico, mas provocam “confusão” e, por isso, deveriam ser desconsiderados, concentrando-se na língua como fato empírico. Além disso, tal posição reforça a ideia de que existiria algum tipo de linguagem essencialmente constativa e, por isso, verdadeira, capaz de fazer uma referência objetiva à realidade, retomando o problema da afirmação na filosofia.

Também é importante lembrar a releitura da obra de Austin feita por John Searle, que teve como principal consequência o aprisionamento do filósofo como um praticante exemplar da Filosofia Analítica. Segundo \textcite{rajagopalan1996austin}, Searle sustentava que o cerne do ato de fala seria seu “conteúdo proposicional”, ou seja, o valor de verdade do enunciado ou a sua referência. As consequências desse argumento são o enquadramento da obra do inglês em um modelo lógico, do qual ele mesmo fez questão tantas vezes de duvidar \cite{austin1962how}.

Este breve panorama é insuficiente para resgatar a complexidade das ideias de Austin e dos grandes debates e apropriações de sua obra. Contudo, ele é fundamental para compreender como a linguagem opera em diversos contextos e práticas sociais, principalmente o que interessa a esse artigo – infraestruturas técnicas digitais e processos de subjetivação. Na próxima seção, abordaremos essa questão a partir da obra da antropóloga brasileira Cesarino e o uso do conceito performatividade, derivado do debate austiniano, para compreender o uso da linguagem em plataformas digitais.


\section{Linguagem, subjetivação e infraestruturas digitais}\label{sec-conduta}
As transformações subjetivas provocadas por mudanças tecnológicas analisadas por Fanon e Cesarino são, de certa forma, diferentes. Ao invés de aparelhos tecnológicos de transmissão e recepção de rádio, como discute \textcite{fanon2021argelia}, nos últimos anos, foi a popularização e o aumento da capilaridade de aparelhos eletrônicos e de plataformas digitais entre a população que provocou uma série de transformações na esfera pública relacionadas à digitalização da vida. Essas mudanças foram impulsionadas, sobretudo, por grandes empresas do setor, conhecidas como Big Techs, que, em sua maioria, estão localizadas no  Vale do Silício nos Estados Unidos, como a Google e a Meta.

Nesse ponto, é preciso apontar uma diferença estrutural dos dois contextos. O rádio da Revolução Argelina funcionava para a formação de comunidades imaginadas e homogêneas, seja de colonos ou de revolucionários, e era baseada em uma lógica de broadcasting de escuta simultânea. Por sua vez, as infraestruturas digitais das quais tratou Cesarino operam com foco no indivíduo, utilizando estratégias de micro-targeting, clusterização e personalização de feeds para a indução de comportamento a partir de dados comportamentais. Assim, enquanto o rádio unia as pessoas em comunidades, as plataformas digitais isolam e invisibilizam a diferença.  Essa diferença, no entanto, não impede nossa visita à argumentação de Fanon para pensar o digital, um esforço reflexivo urgente como apontam \textcite{faustino2023colonialismo}.

Atualmente, a produção de hardwares e softwares ao redor do mundo impulsionou novas cadeias de produção e de consumo, levando a uma nova etapa do capitalismo, caracterizada por tecnologias avançadas, internet das coisas e sistemas algorítmicos. Assim, a informação e os dados passaram a ser recursos cruciais para a produção de mais valia e, por isso, sua criação, monitoramento e exploração comercial tornaram-se fundamentais para o mercado da informação. É consenso entre estudiosas e estudiosos de diferentes campos que tal dinâmica exerce influência, cada vez mais, na criação de infraestruturas técnicas digitais que buscam manter os sujeitos cada vez mais conectados via manipulação algorítmica de afetos \cite{cesarino2021b,faustino2023colonialismo,han2023psicopolitica}.

\textcite[p. 154]{faustino2023colonialismo} sintetizam a modulação do comportamento via exploração, cruzamento e processamento de dados de “capitalismo cognitivo”. Os autores discutem as formas de subjetividade emergentes nesse modo de produção, articulando as transformações objetivas e subjetivas: “A reconfiguração cibernética do corpo, da experiência e da subjetividade se expressa tanto nas alterações objetivas do tempo de circulação de mercadorias, valores e dados quanto na percepção sensitiva a respeito de nós, dos outros e do mundo.” \cite[p. 155]{faustino2023colonialismo}.

\textcite{han2023psicopolitica}, pensador também citado por \textcite{faustino2023colonialismo}, compara os modelos de vigilância do poder dos séculos XIX (o panóptico de Bentham) e XX (o Grande Irmão) com as condições da vigilância digital deste início de século XXI. Ele chama este novo modelo de panóptico digital e identifica no micro-targeting sua ferramenta de subjetivação mais eficiente: “O micro-targeting é aplicado para abordar os eleitores com mensagens direcionadas e personalizadas, e assim influenciá-los. O micro-targeting, como prática da microfísica do poder, é uma psicopolítica movida por dados.” \cite[p. 72]{han2023psicopolitica}.

As consequências dessas mudanças podem ser observadas no processo de digitalização da política, visível nas discussões recentes de Cesarino sobre os públicos antiestruturais no Brasil \cite{cesarino2007reencontro,cesarino2019identidade,cesarino2020populismo,cesarino2021a,cesarino2021b,cesarino2022a,cesarino2022b}. Para ela, o processo de polarização política está relacionado a uma lógica de subjetivação que tem sido chamado de eu-pistemologia, na qual o mundo é cindido entre as categorias “amigo” e “inimigo”, que se sobrepõem aos sistemas de validação social tradicionais institucionais, como a ciência, a democracia e a justiça. Tal processo seria provocado pela formação de públicos clusterizados e pela criação de feeds personalizados, que mobilizam afetos subjetivos para manter os usuários conectados.

Segundo \textcite{cesarino2022b}, o desenho das novas tecnologias e suas infraestruturas caminham para igualar o pensamento animal-humano-máquina por meio da algoritmização da cognição, provocando um processo de reprimarização do cérebro, no qual somos levados a criar relações e identidades via afetos e não por relações lógicas. Por isso, crescem os conteúdos e os públicos antiestruturais ou conspiratórios, como antivacinas e antidemocráticos, impulsionados pela disseminação de notícias falsas e canais de desinformação. Dessa forma, pouco tem tido efeito o apelo à “verdade”, à ciência e às instituições públicas, que nessa dicotomia são vistas como inimigas por muitos e pouco têm a confiança dos que se fecham sobre a própria “verdade”.

Embora não sejam seu foco principal, o performativo e a performatividade ajudam a compreender o caos da política digital. Segundo \textcite[p. 86]{cesarino2021b}, entre os eventos transversais dos desequilíbrios causados pelas plataformas digitais na esfera pública, destaca-se o que ela nomeia como linguagem performativa, que, em suas palavras, é “redundante e reducionista, com fortes componentes estéticos e afetivos. São linguagens que não descrevem uma realidade preexistente, mas geram efeitos sobre, e coproduzem, os sujeitos que comunicam”. Em sua perspectiva, a “performatividade essencialmente antagonística que encontra respaldo na homofilia algorítmica das plataformas”, quando aplicada em processos de neoliberalização e de radicalização, “tornam-se centrais a formas de subjetivação altamente performativas” \cite[p. 86]{cesarino2021b}.

Essa visão da linguagem remete ao performativo de \textcite{austin1962how}, um tipo de enunciado que realiza uma ação ao invés de simplesmente descrevê-la. Os termos “linguagem performativa” ou “performatividade” nunca foram usados por Austin, mas são uma derivação direta da obra do autor, especialmente popularizada nas Ciências Humanas pela interpretação que a filósofa \textcite{butler1997excitable,butler1998fundamentos,butler1999gender,butler2015quadros,butler2017vida} tem feito dessa visão desde os anos 1990.

Contudo, se aprofundarmos nossa leitura da maneira como a performatividade tem sido apropriada na discussão sobre as plataformas digitais, é possível notar que este uso guarda em si um pressuposto que foi duramente combatido pelo próprio \textcite{austin1962how,austin1998performativo}: a suposição de que existiria, em algum outro ponto, uma linguagem constativa, ou seja, verdadeira, factual, lógica e racional, em oposição ao falso, ao redundante, ao reducionismo e aos afetos.

Em uma obra recente, o livro O mundo do avesso: política e verdade na era digital \cite{cesarino2022b}, a autora, ao explicar a materialidade dos sistemas cibernéticos, faz referência a Austin para ressaltar, também, o lado material da linguagem, associado, justamente, à performatividade:

\begin{quote}
    John Austin ([1962] 1990) propôs a noção do performativo para delimitar os enunciados que não descrevem \textbf{o mundo tal como ele já se apresenta}, mas produzem realidades no ato mesmo de serem expressos. Seu valor de verdade se liga ao contexto da enunciação, onde os elementos certos devem estar presentes para que ele seja eficaz \cite[p. 32, grifos nossos]{cesarino2022b}.
\end{quote}

Ao atribuir a noção de valor de verdade ao contexto da enunciação, Cesarino atribui um valor de verdade ao performativo, caso as condições de felicidade do enunciado sejam cumpridas. Contrariamente, na discussão de Austin, a referência à verdade e à falsidade, categorias centrais da Filosofia da Lógica, deixou de ter um papel relevante e foi abandonada pelo próprio filósofo, uma vez que a distinção entre performativos e constativos se mostrou improdutiva \cite{austin1998performativo}. Por isso, ao analisar o exemplo do casamento – proposto pelo filósofo e retomado pela antropóloga –, torna-se menos interessante compreender, para uma visão performativa da linguagem, se o ato proferido pelo padre é verdadeiro ou falso, em vez de quais efeitos foram produzidos por aquele ato de fala.

Embora sejam poucas as referências diretas a Austin na obra de Cesarino, a alusão à performatividade é recorrente, principalmente em textos que abordam o populismo digital e as transformações na esfera pública decorrentes da capilarização das plataformas digitais entre a população. Nessas publicações, são atribuídas características aos discursos performativos associadas à falsidade no sentido lógico, principalmente quando o foco é caracterizar os populistas: “Daí o caráter impreciso, redundante, simplificador, emotivo, ‘vazio’ – em uma palavra, performativo – do discurso populista” \cite[p. 99]{cesarino2020populismo}. Ou ainda:

\begin{quote}
    embora tenham trajetórias historicamente diferentes, o discurso populista e a memética digital apresentam hoje grande convergência em termos de sua linguagem essencialmente performativa, redundante e reducionista, com fortes componentes estéticos afetivos. São linguagens que \textbf{não descrevem uma realidade preexistente}, mas geram efeitos sobre, e coproduzem, os sujeitos que comunicam \cite[p. 86, grifos nossos]{cesarino2021b}.
\end{quote}

Assim, a performatividade aparece associada a duas formas de linguagem específicas, o discurso político de líderes populistas e os memes que circulam em grupos e em plataformas digitais de seus seguidores, ambas consideradas falseadoras de “o mundo tal como ele já se apresenta” ou “uma realidade preexistente”. Dessa forma, ela estaria associada a métricas algorítmicas e a processos de neoliberalização que convergiriam em processos de subjetivação “altamente performativos”, em outras palavras, a performatividade estaria associada à produção de novas identidades que surgem da relação entre usuários aglutinados por plataformas digitais e seus antagonistas, sempre construindo subjetividades alienadas do “mundo” ou “realidade”.

É importante destacar que as referências à performatividade em Cesarino muitas vezes estão associadas à formulação teórica do populismo de \textcite{laclau2005populist}. Por isso, ela é caracterizada como um significante vazio para atender a uma necessidade estrutural de vagueza para a formação de identidades e grupos a partir de uma lógica antecipatória e retroativa: induzindo a coletividade que ele aparenta nomear. Contudo, a partir da crítica de \textcite{butler2024laclau}, é possível afirmar que Laclau se mantém em um âmbito estritamente linguístico, como Benveniste, pois atribui exclusivamente ao nome/significante a força performativa, o que contraria o pressuposto de Austin de que o ato de fala não é redutível a formas proposicionais.

A crítica de Butler reside no fato de que Laclau restringe a performatividade política primordialmente ao domínio da análise retórica e da representação linguística, conferindo ao nome (o significante) o poder quase exclusivo de unificar e trazer um sujeito político à existência de forma retroativa e antecipatória. Butler argumenta que essa perspectiva falha ao não reconhecer que a performatividade é também corporal e gestual, não se reduzindo a afirmações proposicionais ou à simples nomeação de uma exigência. Para ela, a contiguidade corpórea em manifestações e assembleias constitui uma forma de significação política autônoma, na qual os corpos, ao aparecerem publicamente, já reivindicam o direito de existir e florescer, evidenciando que os sujeitos são agidos por normas e histórias sociais antes mesmo de poderem agir ou serem nomeados dentro de uma cadeia de equivalência. Além disso, considerar que a subjetivação ocorre apenas no momento da nomeação é ignorar que o sujeito é interpelado por uma matriz de desejo, histórias e relações sociais que ele não escolheu, mas que são cruciais para sua sobrevivência \cite{butler2024laclau}.

Nossa leitura é que tanto Laclau quanto Cesarino acabam reforçando a ideia de uma relação direta entre significante e significado, enfraquecendo o potencial crítico da performatividade em relação a normatividades impostas por infraestruturas digitais.

Mais uma vez, é preciso questionar: se tais linguagens redundantes e reducionistas – que se referem a públicos antiestruturais, como grupos conspiratórios, extrema direita, antivacinas – são caracterizadas como performativas, presume-se que há outros discursos caracterizados como constativos que poderiam evidenciar de forma objetiva uma verdade, “um mundo tal qual ele é”? Estaria o debate sobre manipulação algorítmica evocando um conteúdo proposicional emergido de uma realidade prévia, um contexto original, a partir do qual se estabeleceria a solução contra essa psicopolítica dos dados \cite{han2023psicopolitica}, uma linguagem supostamente lógica, racional e objetiva?

Nesse sentido, na tentativa de responder a essas questões, alguns trabalhos atribuem centralidade à noção de contexto, como faz a própria \textcite{cesarino2021b} para discutir o problema da “pós-verdade”. Atualmente, as condições que determinam o contexto são radicalmente diferentes das propostas por Austin para discutir as condições de felicidade do ato de fala. Isso leva muitos a duvidarem da possibilidade de determiná-lo nas sociedades digitais. Daí o sucesso da noção de colapso de contexto, proposta por \textcite{marwick2010tweet}. As autoras apresentam essa noção para dar conta da multiplicação e sobreposição de audiências em uma única timeline de redes digitais. Nessa perspectiva, nas interações digitais, há uma indeterminação da audiência e um decorrente engajamento em negociações complexas de identidade, impressões e face \cite[p. 123]{marwick2010tweet}. Essas características do contexto digital colapsado afetam as noções de público e privado, levam à intensificação do monitoramento da autoimagem e complexificam as noções de tempo e espaço.

Entre essas características, notamos um destaque para a perda de delimitações contextuais que eram tradicionalmente consideradas importantes para regular as trocas simbólicas pela linguagem. Essa perda faria com que um enunciado possa ser continuamente tirado de um contexto e recontextualizado em novos enquadramentos que evocam novos efeitos e significados, eventualmente até mesmo opostos.

No entanto, depois de várias décadas de crítica ao modelo saussureano de signo, já sabemos que essa potencialidade de recontextualização de qualquer signo não é uma característica emergente unicamente no ambiente digital. O significado da interação nunca é fixo nem está presente de forma imediata – “a presença, o privilégio de sua proximidade com o logos” \cite[p. 31]{derrida1967grammatologie} – ou delimitado pelo contexto; pelo contrário, sua determinação é sempre adiada, “um restante não presente de uma marca diferencial cortada da sua suposta ‘produção’ ou origem” \cite[p. 32]{derrida1990limited}, pois ele depende de um sistema de diferenças definidas por enquadramentos nunca estáveis e sempre em disputa.

Os eventos interacionais digitais rompem invariavelmente consigo mesmos e suas delimitações, numa escala mais rápida do que a interação face-a-face, movendo-se através da rede em cadeias de significação na qual o próprio movimento depende do rompimento com cada contexto anterior. Assim, os atos de fala e o contexto se comportam como os enquadramentos discutidos por \textcite[p. 26]{butler2015quadros}:

\begin{quote}
    O enquadramento que busca conter, transmitir e determinar o que é visto (e algumas vezes, durante um período, consegue fazer exatamente isso) depende das condições de reprodutibilidade para ter êxito. Essa própria reprodutibilidade, porém, demanda uma constante ruptura com o contexto, única constante delimitação de novos contextos, o que significa que o “enquadramento” não é capaz de conter completamente o que transmite e se rompe toda vez que tenta dar uma organização definitiva a seu conteúdo. Em outras palavras, o enquadramento não mantém nada integralmente em um lugar, mas ele mesmo se torna uma espécie de rompimento perpétuo, sujeito a uma lógica temporal de acordo com a qual se desloca de um lugar para outro. Como o enquadramento rompe constantemente com seu contexto, esse autorrompimento converte-se em parte de sua própria definição.
\end{quote}

Aquilo que a sociedade ocidental insistiu em chamar de “verdade” como reflexo de uma “realidade preexistente”, nos ambientes digitais, enfrenta o risco estrutural do funcionamento do signo e seu enquadramento, desde sempre autorrompido. Os contextos são colapsados pelas condições escalares de reprodutibilidade dos atos de fala em ambientes digitais. Somada a isso, a eu-pistemologia impõe uma barreira para uma suposta verdade universal, institucionalizada por sistemas de peritos como o judiciário, a ciência ou a democracia. É a “verdade” individual internalizada via afetos direcionados por algoritmos em plataformas digitais, articulada ao sentimento do sujeito de “estar no controle quando, na verdade, acaba interpelado por estruturas que mal conhece e controla, como é o caso do big data e seus algoritmos frios” \cite[p. 156]{faustino2023colonialismo}.

Nesse cenário, reiterando a posição de Austin, a busca por uma “verdade” original ou lógica se torna improdutiva, haja visto o enorme volume de dados e reenquadramentos que podem transformar qualquer contexto. Dessa forma, a oposição dos valores verdadeiro-falso é pouco relevante para a compreensão de como a linguagem funciona em ambientes digitais, e ainda menos para a ação política frente à produção da desinformação e da ignorância.

Para solucionar tais problemas, o foco deve ser nos efeitos produzidos pelos atos de fala, sua capacidade de criar a realidade, posição adotada – mas um pouco esquecida depois – pela própria \textcite{cesarino2007reencontro} ao analisar pronunciamentos presidenciais de Lula antes da popularização do consumo de smartphones:

\begin{quote}
    Como se pode notar, ambos os pronunciamentos visam à produção de efeitos sobre o receptor da mensagem; não se pautam, portanto, por uma lógica referencial - ou, nos termos de Austin (1975), são atos ilocucionários. O principal destes efeitos (perlocucionários) parece ser o estabelecimento de uma contiguidade e a construção de \textit{um sujeito comum} a partir de emissor e receptor. Além disso, e novamente tomando Austin, os discursos são sujeitos a julgamentos normativos de felicidade e infelicidade, e não a testes racionais de verdade e falsidade. Desta forma, ao “falar politicamente” - no sentido da transmissão da mensagem entre representante e representados - “dizer algo” é sempre “fazer algo” \cite[p. 193, grifos da autora]{cesarino2007reencontro}.
\end{quote}

Buscando tomar “o mundo dos signos” no sentido político de Fanon e a radicalidade do performativo como “fazer algo”, na próxima seção, apresentamos as apropriações de Butler e Derrida da obra de Austin como uma elaboração crítica à velha oposição performativo-constativo e sua contraparte lógica verdade-falsidade, para compreender o funcionamento da linguagem em infraestruturas digitais.

\section{Performatividade e subjetivação}\label{sec-fmt-manuscrito}
Como já indicamos, a Teoria dos Atos de Fala proposta pelo filósofo britânico possibilitou o desenvolvimento de diversas perspectivas teóricas, bem como originou um campo de estudos na linguística conhecido como Pragmática, preocupado em investigar a linguagem em uso. Conforme aponta \textcite{pinto2007conexoes,pinto2009corpo}, é possível traçar aproximações e percursos teóricos decorrentes do trabalho de Austin no que concerne ao debate crítico sobre linguagem e suas articulações com subjetivação e identidades, destacando, principalmente, a leitura crítica de \textcite{derrida1990limited} e as apropriações radicais do ato de fala como ato de corpo de \textcite{felman1980scandale,butler1997excitable,butler1999gender}.

Seguindo por esse caminho trilhado por essas autoras e autor, argumentamos que \textcite{austin1962how} partiu da ideia de um sujeito intencional e consciente da totalidade de seu ato de fala, prevendo uma unidade de sentido na sua realização e a força ilocucionária, a princípio, partiria unicamente desse sujeito. Contudo, tal questão foi antecipada pelo próprio Austin com a noção de uptake – um reconhecimento mútuo entre falantes de que algo está assegurado ou ritualizado por meio do ato de fala, pois não haveria critérios formais definitivos de sucesso desse ato \textcite[p. 120–146]{austin1962how}.

\textcite[p. 28]{derrida1990limited} interpreta essa ausência de critérios formais como a própria condição para legibilidade dos signos, “a disruptura, em última análise, da autoridade do código como sistema finito de regras; a destruição radical, ao mesmo tempo, de todo contexto como protocolo do código”. Essa disruptura é a própria condição de funcionamento do signo, que, para funcionar como signo, deve continuar legível mesmo na ausência ou no desaparecimento de quem o produziu, ou mesmo na falta de sustentação de uma intenção ou atenção presente ou um querer-dizer \cite[p. 29]{derrida1990limited}. Dessa forma,

\begin{quote}
    O uptake desfaz a possibilidade de “falante consciente da totalidade do ato” porque exige alteridade, descentraliza o falante, fragmenta assim os sentidos e os efeitos, deixando portanto escapar “restos”, produzindo uma “polissemia irredutível” (Derrida 1990:39) própria à performatividade \cite[p. 8]{pinto2007conexoes}.
\end{quote}

Em torno desse debate, Derrida formula a noção de iterabilidade, ou seja, a possibilidade estrutural de todo signo ser repetido na ausência de seu referente, de seu significado e de sua intenção determinada \cite{derrida1990limited,pinto2009corpo}. Sob essa lógica, o rito se torna essencial para os atos de fala, uma vez que sua iterabilidade – a repetição na diferença – aciona a citacionalidade necessária para o funcionamento do signo. O próprio Austin suspeita das determinantes contextuais e prefere a força como categoria de análise do funcionamento do ato de fala, porque significado (o centro da explicação contextual) equivale a sentido e referência, categorias usadas na Filosofia da Lógica para analisar o valor de verdade de uma sentença:

\begin{quote}
    [...] há alguns anos, temos percebido cada vez mais claramente que a ocasião de um enunciado importa seriamente e que as palavras usadas devem, até certo ponto, ser “explicadas” pelo “contexto” em que foram concebidas para serem ou foram efetivamente proferidas em um intercâmbio linguístico. No entanto, talvez ainda sejamos muito propensos a dar essas explicações em termos dos “significados das palavras”. É verdade que podemos usar “significado” também com referência à força ilocucionária — “Ele quis dizer isso como uma ordem” etc. Mas quero distinguir força e significado no sentido em que significado é equivalente a sentido e referência, assim como se tornou essencial distinguir sentido e referência dentro do significado \cite[p. 100]{austin1962how}.
\end{quote}

Em sua interpretação de Austin, \textcite[p. 29]{derrida1990limited} argumenta que o contexto está submetido a uma indeterminação e à impossibilidade ou insuficiência de sua saturação empírica. Essa necessidade de ritualização implica uma articulação entre determinadas circunstâncias no tempo, fazendo com que todo ato ilocucionário exceda sua temporalidade imediata e impedindo, desde seu início, que possamos atingir a totalidade da situação desse ato de fala \cite{butler1997excitable}. Assim, a força do ato de fala é atualizada no momento de seu acontecimento, fortalecido pela força do rito, mostrando-se como um efeito histórico \cite{butler1997excitable} e, portanto, evidenciando que o performativo é um acontecimento que não espera a deliberação, a consciência e a organização do sujeito \cite{pinto2009corpo}.

A leitura crítica de Derrida e a apropriação de \textcite{felman1980scandale} inspiraram \textcite{butler1997excitable} a teorizar o corpo a partir dos atos de fala. Sua principal inovação está em incluir o ponto-cego da fala, o corpo, na análise, reforçando, assim, a relação entre o linguístico e o político \cite{pinto2009corpo}. Para Butler, a interpelação do sujeito pelo ato de fala marca a inscrição da cultura na materialidade corporal e, consequentemente, leva à sujeição a sistemas coercitivos. Por isso, uma de suas preocupações é a criação de condições linguísticas de sobrevivência que possam operar e agir nos termos do poder. Nessa perspectiva, Butler aponta uma estrutura ambivalente de performatividade que regula a subjetivação.

Tal ideia de ambivalência da subjetivação se relaciona com a forma como somos constituídos por um poder exterior que nos submete e determina nossas condições de existência \cite{butler2017vida}. Contudo, esse poder é também produtivo na medida em que, por meio dele, constituímo-nos como sujeitos. Dessa forma, a subjetividade é formada por um processo ambivalente entre sujeição e subjetivação. Em \textcite{butler1997excitable}, essa argumentação encontra ecos no debate sobre linguagem e suas características performativas, pois, em nossa constituição como sujeitos, somos submetidos a ela, logo, o corpo – a materialidade produzida e produtora da linguagem – torna-se um feito ameaçado pela própria linguagem. Além disso, é ela que constitui seus efeitos possíveis, evocando sedimentações históricas, determinando espaços de inteligibilidade, de regulação e de legitimação. Ainda assim, a ambivalência permanece, afinal, sempre há a possibilidade de escapar a determinismos interrompendo a cadeia de citações dos signos, provocando pontos de deslizamentos, de descontinuações e de rearticulação, o que leva, consequentemente, a uma operação conjunta entre a violência coercitiva e a produtiva \cite{pinto2009corpo}.

Neste ponto, é possível articular esse debate a uma teoria da linguagem em ambientes virtuais para além da distinção entre “linguagens performativas” e “constativas”, entre “verdade” e “falsidade”. As posições adotadas por Butler e Derrida pouco se interessam por essas distinções e focam, sobretudo, nos efeitos produzidos pela linguagem. Dessa forma, não haveria uma “linguagem constativa” capaz de acessar a “verdade” preexistente ao ato de fala, pois todas as formas de linguagem seriam performativas. Assim, a linguagem utilizada no populismo digital e nos memes associados a ele não seria performativa porque é reducionista. Na concepção forte de performatividade, mesmo a linguagem utilizada no sistema de peritos tradicionais, como na mídia tradicional, nos partidos políticos e na ciência, é, também, performativa.

Considerando que muitas vezes a internet é entendida apenas como um “contexto” específico de interação e muito pouco confrontada com as especificidades de suas infraestruturas e condições de reprodutibilidade, é essencial pensar o digital como uma psicopolítica, mas um pouco além de uma microfísica do poder \cite{han2023psicopolitica}. É preciso encarar esses processos subjetivantes em relação aos excessos, ou seja, encarar a relação dialética entre a capacidade de agir imbricada nas formas históricas e alteração da subjetividade humana diante deste “salto tecnológico” \cite[p. 153]{faustino2023colonialismo}.

Diante das alterações diárias que as infraestruturas digitais impõem às nossas vidas, temos a sensação de que algo maior está em andamento. Há ainda espaço para “uma outra verdade concreta” \cite{fanon2021argelia}, uma “verdade” não constativa, mas constituída pela força de um performativo que torna possível contrapontos a essa psicopolítica?

Para \textcite{fanon2021argelia}, apenas a “apropriação anticolonial de algumas tecnologias sociais” \cite[p. 186]{faustino2023colonialismo} pode promover uma “dialética da dominação e da libertação através da descolonização da tecnologia de comunicação” \cite[p. 187]{faustino2023colonialismo}. O que isso pode nos demandar diante das infraestruturas digitais ainda não está resolvido. Só um processo crítico revolucionário pode disputar a liberdade transformada em subproduto da psicopolítica digital \cite[p. 72–73]{han2023psicopolitica} e contrapor outras “verdades”. Seja como for, é certamente a performatividade da linguagem como jogo de forças, atos de criação de realidades e de capacidades de agir, que vai mover a formação de novas vozes anti-imperialistas diante das infraestruturas digitais.


\section{Conclusão}\label{sec-formato}
Neste artigo, buscamos explorar como a performatividade ajudaria a compreender os impactos das infraestruturas digitais no mundo contemporâneo, principalmente no que se refere à articulação entre linguagem e formas de subjetivação mediadas pela tecnologia. Para isso, partimos da análise de Fanon sobre o uso do rádio durante a Revolução Argelina, afinal, o autor evidencia como as transformações tecnológicas implicam transformações subjetivas. A contraposição de verdades entre colonizadores e colonizados sobre a realidade colonial evidencia a força performativa da linguagem: criar, realizar, transformar realidades e sujeitos com palavras.

Embora o psiquiatra martinicano não tenha mencionado o termo performativo, em nossa leitura, sua análise sobre a Revolução ecoa a Teoria dos Atos de Fala de Austin, uma vez que, por meio dela, podemos retomar a crítica do filósofo britânico sobre a questão da objetividade e da verdade na filosofia analítica. Contrapondo-nos às vertentes que, de certa forma, aprisionam a teoria de Austin em termos que o próprio descartou, como “verdadeiro” e “falso” ou proposições linguísticas, acreditamos que a linguagem não é meramente descritiva da realidade ou que funcione de forma objetiva. Assim como na oposição do discurso colonizador ao colonizado, a linguagem cria as realidades que aparenta descrever. Nesse arcabouço, ao invés de pensar a verdade como categoria analítica para explicar textos anti-factuais em circulação, seguimos Fanon e Austin e lemos verdade como um desdobramento da luta pela linguagem, um efeito perlocucionário.

Além disso, a categoria do performativo também é observada nos trabalhos de Cesarino sobre o populismo digital e a ascensão de públicos antiestruturais no Brasil. Contudo, ao evocar a performatividade para compreender a formação de tais grupos, ela nos coloca diante de uma escolha epistemológica e política: devemos voltar à questão da verdade, buscando restaurar o status da linguagem constativa, portanto “verdadeira” ou “falsa”; ou devemos abandoná-la – como fizeram Austin, Butler e Derrida – e nos preocuparmos com os efeitos de subjetivação produzidos pela força e pelo risco dos performativos?

Em nossa perspectiva, compreender as transformações subjetivas catalisadas por infraestruturas digitais exige uma postura crítica em relação a essa questão. A dialética entre capacidade de agir e as alterações na subjetividade provocadas pelas tecnologias da comunicação exige uma análise focada nos efeitos produzidos pela linguagem, haja vista que a contraposição de verdades pouco tem efetividade crítica num contexto marcado por sistemas de referência focados no indivíduo e em seus afetos. Por isso, interrogamos a obra de Cesarino não por um purismo teórico, mas como uma ampliação crítica focada no resgate do potencial e da força do performativo e da performatividade para a compreensão dos dilemas da linguagem numa sociedade digital.

O pressuposto de que toda linguagem seja essencialmente performativa é fundamental para a compreensão de seus efeitos materiais e normativos, que constituem os sujeitos e suas condições de ação. Dessa forma, nas plataformas digitais, a performatividade adquire uma dimensão tecnopoliticamente mediada. Os atos de fala não se realizam apenas por regras sociais compartilhadas, mas também em conformidade com as arquiteturas algorítmicas que regulam a circulação dos discursos. A subjetividade, nesse ambiente, não é apenas constituída pela linguagem, mas também moldada por sistemas automáticos de curadoria, moderação e ranqueamento que operam segundo lógicas de engajamento, atenção e lucro.

Nossa leitura indica que a ação subjetiva nas redes é ambígua e tensionada. Por um lado, há uma forte interpelação algorítmica que limita as possibilidades de fala e de visibilidade; por outro, há brechas performativas – espaços de reinscrição, crítica, resistência e reapropriação – que podem ser mobilizadas estrategicamente pelos sujeitos. A capacidade de agir, portanto, não é eliminada, mas reconfigurada. Ela exige um novo entendimento da autonomia, não mais como independência plena, mas como agência situada, exercida em meio a normas, códigos e dispositivos técnicos.

Construímos nossa argumentação para avançar mais o debate sobre a performatividade como categoria para a compreensão da capacidade de agir pela linguagem em infraestruturas digitais. Dados os limites de extensão do artigo, esperamos que este primeiro diálogo estimule estudos empíricos que enfrentem os problemas apontados.

Concluímos, até este momento, que a noção de performatividade continua sendo uma ferramenta teórica poderosa para compreender os processos de subjetivação na era digital, desde que seja pensada criticamente e em articulação com as materialidades técnicas que a condicionam. Ao enfatizar os efeitos da linguagem e os modos como sujeitos são constituídos nas e pelas plataformas, reiteramos a importância de compreender a ação não como livre escolha ou como submissão total, mas como um campo de disputa onde operam simultaneamente determinação estrutural e potência de transformação.


\printbibliography\label{sec-bib}
% if the text is not in Portuguese, it might be necessary to use the code below instead to print the correct ABNT abbreviations [s.n.], [s.l.]
%\begin{portuguese}
%\printbibliography[title={Bibliography}]
%\end{portuguese}


%full list: conceptualization,datacuration,formalanalysis,funding,investigation,methodology,projadm,resources,software,supervision,validation,visualization,writing,review
\begin{contributors}[sec-contributors]
\authorcontribution{Hugo de Oliveira}[conceptualization,datacuration,formalanalysis,investigation,methodology,projadm,writing,review]
\authorcontribution{Joana Plaza Pinto}[conceptualization,formalanalysis,investigation,methodology,projadm,supervision,visualization,writing,review]
\end{contributors}


\begin{dataavailability}
\txtdataavailability{databody} % options: dataavailable, dataonly, databody, datanotav, nodata
\end{dataavailability}




\end{document}

